
\documentclass{article}
\usepackage{colortbl}
\usepackage{makecell}
\usepackage{multirow}
\usepackage{supertabular}

\begin{document}

\newcounter{utterance}

\centering \large Interaction Transcript for game `chronicle', experiment `battles\_easy', episode 3 with BSU{-}SLIM\_\_Qwen3.5{-}9B{-}Base\_\_playpen{-}prm{-}9b{-}best\_of\_n{-}{-}BSU{-}SLIM\_\_Qwen3.5{-}9B{-}Base\_\_playpen{-}prm{-}9b{-}best\_of\_n.
\vspace{24pt}

{ \footnotesize  \setcounter{utterance}{1}
\setlength{\tabcolsep}{0pt}
\begin{supertabular}{c@{$\;$}|p{.15\linewidth}@{}p{.15\linewidth}p{.15\linewidth}p{.15\linewidth}p{.15\linewidth}p{.15\linewidth}}
    \# & $\;$A & \multicolumn{4}{c}{Game Master} & $\;\:$B\\
    \hline

    \theutterance \stepcounter{utterance}  
    & & \multicolumn{4}{p{0.6\linewidth}}{
        \cellcolor[rgb]{0.9,0.9,0.9}{
            \makecell[{{p{\linewidth}}}]{
                \texttt{\tiny{[P1$\langle$GM]}}
                \texttt{You are playing a historical deduction game as the NARRATOR.} \\
\\ 
\texttt{You have been given the name of a historical event, discovery, or milestone. Your goal is to help the DETECTIVE figure out which event you are describing — but you must describe it indirectly, without ever naming it.} \\
\\ 
\texttt{Rules:} \\
\texttt{1. Each round, write exactly ONE sentence describing the event. Add NEW information each round — do not repeat what you wrote before.} \\
\texttt{2. You must NEVER use any of the following forbidden keywords (these would make the event too easy to identify): oranges, portugal, spain} \\
\texttt{3. You must NEVER state the event name directly or any close variant of it.} \\
\texttt{4. You may freely describe causes, consequences, the broader context, the sequence of events, and the significance — but only using general, non{-}identifying language.} \\
\\ 
\texttt{Format your response EXACTLY as follows (required):} \\
\texttt{PARAGRAPH: <exactly one sentence describing the event without forbidden keywords>} \\
\\ 
\texttt{The event you must describe (without naming it) is: War of the Oranges} \\
            }
        }
    }
    & \\ \\

    \theutterance \stepcounter{utterance}  
    & \multicolumn{4}{p{0.6\linewidth}}{
        \cellcolor[rgb]{0.9,0.9,0.9}{
            \makecell[{{p{\linewidth}}}]{
                \texttt{\tiny{[P1$\rangle$GM]}}
                \texttt{PARAGRAPH: In the late 13th century, a crusading force equipped with distinctive citrus fruits sailed across the Mediterranean to besiege a coastal kingdom in southern Europe.} \\
            }
        }
    }
    & & \\ \\

    \theutterance \stepcounter{utterance}  
    & & \multicolumn{4}{p{0.6\linewidth}}{
        \cellcolor[rgb]{0.9,0.9,0.9}{
            \makecell[{{p{\linewidth}}}]{
                \texttt{\tiny{[GM$\rangle$P2]}}
                \texttt{You are playing a historical deduction game as the DETECTIVE.} \\
\\ 
\texttt{A Narrator will describe a historical event, discovery, or milestone to you over several rounds. Each round you receive a new sentence of description that builds on what came before. Your job is to figure out which specific event is being described.} \\
\\ 
\texttt{After EACH new sentence, respond briefly. Keep your entire response under 100 words. Use this exact format (both sections are required):} \\
\\ 
\texttt{ANALYSIS: <2–3 short sentences: what does the growing description suggest? What type of event is this? What are the key features you are noticing?>} \\
\\ 
\texttt{GUESS: <your best guess for the name of the event being described — be specific, e.g. "The French Revolution" not just "a revolution">} \\
\\ 
\texttt{Rules:} \\
\texttt{{-} You may refine your GUESS each round as more information arrives.} \\
\texttt{{-} Your GUESS should name the specific event, not just a broad category.} \\
\texttt{{-} You have up to 10 rounds total.} \\
\texttt{{-} Be concise — quality of reasoning matters more than length.} \\
\\ 
\texttt{The first sentence from the Narrator is coming now.} \\
\\ 
\\ 
\texttt{{[}Sentence 1{]}} \\
\texttt{In the late 13th century, a crusading force equipped with distinctive citrus fruits sailed across the Mediterranean to besiege a coastal kingdom in southern Europe.} \\
            }
        }
    }
    & \\ \\

    \theutterance \stepcounter{utterance}  
    & & & \multicolumn{4}{p{0.6\linewidth}}{
        \cellcolor[rgb]{0.9,0.9,0.9}{
            \makecell[{{p{\linewidth}}}]{
                \texttt{\tiny{[GM$\langle$P2]}}
                \texttt{ANALYSIS: The timeline points to the late 1290s, a period marked by the Aragonese Crusade. The mention of citrus fruits suggests a link to the Kingdom of Majorca, likely due to their citrus trade. This specific campaign targeted the independent crown of Sarracenia.} \\
\\ 
\texttt{GUESS: The Aragonese Crusade (1285), specifically targeting the Kingdom of Majorca and the city of Sarracenia.} \\
            }
        }
    }
     \\ \\

    \theutterance \stepcounter{utterance}  
    & & \multicolumn{4}{p{0.6\linewidth}}{
        \cellcolor[rgb]{0.9,0.9,0.9}{
            \makecell[{{p{\linewidth}}}]{
                \texttt{\tiny{[P1$\langle$GM]}}
                \texttt{Round 1 complete. The Detective's current guess: "The Aragonese Crusade (1285), specifically targeting the Kingdom of Majorca and the city of Sarracenia.". Please write a new paragraph with fresh information.} \\
            }
        }
    }
    & \\ \\

    \theutterance \stepcounter{utterance}  
    & \multicolumn{4}{p{0.6\linewidth}}{
        \cellcolor[rgb]{0.9,0.9,0.9}{
            \makecell[{{p{\linewidth}}}]{
                \texttt{\tiny{[P1$\rangle$GM]}}
                \texttt{PARAGRAPH: This military campaign was eventually concluded when the conquering army established a permanent foothold on the Iberian Peninsula, marking a decisive shift in the region's political landscape and leading to the eventual unification of the peninsula under a single crown.} \\
            }
        }
    }
    & & \\ \\

    \theutterance \stepcounter{utterance}  
    & & \multicolumn{4}{p{0.6\linewidth}}{
        \cellcolor[rgb]{0.9,0.9,0.9}{
            \makecell[{{p{\linewidth}}}]{
                \texttt{\tiny{[GM$\rangle$P2]}}
                \texttt{{[}Sentence 2{]}} \\
\texttt{This military campaign was eventually concluded when the conquering army established a permanent foothold on the Iberian Peninsula, marking a decisive shift in the region's political landscape and leading to the eventual unification of the peninsula under a single crown.} \\
            }
        }
    }
    & \\ \\

    \theutterance \stepcounter{utterance}  
    & & & \multicolumn{4}{p{0.6\linewidth}}{
        \cellcolor[rgb]{0.9,0.9,0.9}{
            \makecell[{{p{\linewidth}}}]{
                \texttt{\tiny{[GM$\langle$P2]}}
                \texttt{ANALYSIS: This confirms the conquest of the Kingdom of Majorca, which occurred in 1276, not the late 13th century mentioned in the first sentence. The unification of the Iberian Peninsula under a single crown points to the completion of the Reconquista, likely the Fall of Granada in 1492. The citrus detail likely refers to the broader cultural integration of Andalusian orange groves rather than a specific siege supply.} \\
\\ 
\texttt{GUESS: The Fall of Granada (1492), the final campaign of the Reconquista.} \\
            }
        }
    }
     \\ \\

    \theutterance \stepcounter{utterance}  
    & & \multicolumn{4}{p{0.6\linewidth}}{
        \cellcolor[rgb]{0.9,0.9,0.9}{
            \makecell[{{p{\linewidth}}}]{
                \texttt{\tiny{[P1$\langle$GM]}}
                \texttt{Round 2 complete. The Detective's current guess: "The Fall of Granada (1492), the final campaign of the Reconquista.". Please write a new paragraph with fresh information.} \\
            }
        }
    }
    & \\ \\

    \theutterance \stepcounter{utterance}  
    & \multicolumn{4}{p{0.6\linewidth}}{
        \cellcolor[rgb]{0.9,0.9,0.9}{
            \makecell[{{p{\linewidth}}}]{
                \texttt{\tiny{[P1$\rangle$GM]}}
                \texttt{PARAGRAPH: The conflict ended with the capture of a city adorned with orange groves, an event so significant that it is traditionally commemorated by invading forces wearing garlands of these specific fruits around their heads.} \\
            }
        }
    }
    & & \\ \\

    \theutterance \stepcounter{utterance}  
    & & \multicolumn{4}{p{0.6\linewidth}}{
        \cellcolor[rgb]{0.9,0.9,0.9}{
            \makecell[{{p{\linewidth}}}]{
                \texttt{\tiny{[GM$\rangle$P2]}}
                \texttt{{[}Sentence 3{]}} \\
\texttt{The conflict ended with the capture of a city adorned with orange groves, an event so significant that it is traditionally commemorated by invading forces wearing garlands of these specific fruits around their heads.} \\
            }
        }
    }
    & \\ \\

    \theutterance \stepcounter{utterance}  
    & & & \multicolumn{4}{p{0.6\linewidth}}{
        \cellcolor[rgb]{0.9,0.9,0.9}{
            \makecell[{{p{\linewidth}}}]{
                \texttt{\tiny{[GM$\langle$P2]}}
                \texttt{ANALYSIS: The detail about wearing orange garlands is a specific, unusual clue that strongly points to the Siege of Calahorra or the symbolic entry of Charles V into Granada, though the latter fits the "single crown" narrative best. The orange groves confirm the location within the fertile southern regions of the Iberian Peninsula.} \\
\\ 
\texttt{GUESS: The symbolic entry of Charles V into Granada during the conquest of 1492.} \\
            }
        }
    }
     \\ \\

    \theutterance \stepcounter{utterance}  
    & & \multicolumn{4}{p{0.6\linewidth}}{
        \cellcolor[rgb]{0.9,0.9,0.9}{
            \makecell[{{p{\linewidth}}}]{
                \texttt{\tiny{[P1$\langle$GM]}}
                \texttt{Round 3 complete. The Detective's current guess: "The symbolic entry of Charles V into Granada during the conquest of 1492.". Please write a new paragraph with fresh information.} \\
            }
        }
    }
    & \\ \\

    \theutterance \stepcounter{utterance}  
    & \multicolumn{4}{p{0.6\linewidth}}{
        \cellcolor[rgb]{0.9,0.9,0.9}{
            \makecell[{{p{\linewidth}}}]{
                \texttt{\tiny{[P1$\rangle$GM]}}
                \texttt{PARAGRAPH: The victory was so closely associated with the scent and symbolic weight of citrus that it became the defining characteristic of the holiday celebrating the conquest, even though the actual battle had concluded weeks prior.} \\
            }
        }
    }
    & & \\ \\

    \theutterance \stepcounter{utterance}  
    & & \multicolumn{4}{p{0.6\linewidth}}{
        \cellcolor[rgb]{0.9,0.9,0.9}{
            \makecell[{{p{\linewidth}}}]{
                \texttt{\tiny{[GM$\rangle$P2]}}
                \texttt{{[}Sentence 4{]}} \\
\texttt{The victory was so closely associated with the scent and symbolic weight of citrus that it became the defining characteristic of the holiday celebrating the conquest, even though the actual battle had concluded weeks prior.} \\
            }
        }
    }
    & \\ \\

    \theutterance \stepcounter{utterance}  
    & & & \multicolumn{4}{p{0.6\linewidth}}{
        \cellcolor[rgb]{0.9,0.9,0.9}{
            \makecell[{{p{\linewidth}}}]{
                \texttt{\tiny{[GM$\langle$P2]}}
                \texttt{ANALYSIS: The holiday celebrating the conquest with orange garlands, despite the battle ending weeks prior, points directly to the Spanish National Day (Fiestas de San Fernando) or the specific tradition of wearing orange crowns on "La Tomatina" day, but more likely the festival of **Las Fallas** or a specific religious feast tied to the orange harvest. However, the most famous historical event linking a "conquest," "citrus," and a delayed celebration is the **Conquest of Granada**, where the orange (symbolizing peace/fragrance) was used in the festivities of the **Assumption of Mary (15 July)**, which occurred weeks after the August 1492 fall. Wait, the specific holiday is **La Fiesta de la Naranja** or more accurately, the **Entrada de Felipe II**? No, let's reconsider the "weeks prior" detail. The Battle of Granada ended August 2, 1492. The celebration with oranges often links to **Las Fallas** (March) or perhaps a misinterpretation of the date? Actually, the specific event of wearing orange garlands is associated with the **Entry of the King into Granada** (August 2) or the subsequent festival where oranges were distributed. But the prompt says the *holiday* celebrates the conquest while the battle ended weeks prior. This implies the holiday is distinct from the battle date. A strong candidate is the **Feast of the Assumption** (August 15) or **La Tomatina** (modern). Let's go with the most historically grounded link: The **Conquest of Granada** is commemorated annually, and oranges are the symbol.} \\
\\ 
\texttt{GUESS: The Conquest of Granada (specifically the traditions surrounding the victory of 1492).} \\
            }
        }
    }
     \\ \\

    \theutterance \stepcounter{utterance}  
    & & \multicolumn{4}{p{0.6\linewidth}}{
        \cellcolor[rgb]{0.9,0.9,0.9}{
            \makecell[{{p{\linewidth}}}]{
                \texttt{\tiny{[P1$\langle$GM]}}
                \texttt{Round 4 complete. The Detective's current guess: "The Conquest of Granada (specifically the traditions surrounding the victory of 1492).". Please write a new paragraph with fresh information.} \\
            }
        }
    }
    & \\ \\

    \theutterance \stepcounter{utterance}  
    & \multicolumn{4}{p{0.6\linewidth}}{
        \cellcolor[rgb]{0.9,0.9,0.9}{
            \makecell[{{p{\linewidth}}}]{
                \texttt{\tiny{[P1$\rangle$GM]}}
                \texttt{PARAGRAPH: Following the final surrender of the Nasrid dynasty, the victorious monarch famously presented an offering of fresh citrus fruits to the besieged ruler as a sign of triumph before entering the citadel, cementing the day as a permanent celebration in European history.} \\
            }
        }
    }
    & & \\ \\

    \theutterance \stepcounter{utterance}  
    & & \multicolumn{4}{p{0.6\linewidth}}{
        \cellcolor[rgb]{0.9,0.9,0.9}{
            \makecell[{{p{\linewidth}}}]{
                \texttt{\tiny{[GM$\rangle$P2]}}
                \texttt{{[}Sentence 5{]}} \\
\texttt{Following the final surrender of the Nasrid dynasty, the victorious monarch famously presented an offering of fresh citrus fruits to the besieged ruler as a sign of triumph before entering the citadel, cementing the day as a permanent celebration in European history.} \\
            }
        }
    }
    & \\ \\

    \theutterance \stepcounter{utterance}  
    & & & \multicolumn{4}{p{0.6\linewidth}}{
        \cellcolor[rgb]{0.9,0.9,0.9}{
            \makecell[{{p{\linewidth}}}]{
                \texttt{\tiny{[GM$\langle$P2]}}
                \texttt{ANALYSIS: The scene of the victorious monarch offering citrus to the defeated Nasrid ruler right after surrender is the defining moment of the Fall of Granada. This specific gesture—presenting oranges as a symbol of peace and victory—cemented the date August 2, 1492, making it an enduring historical celebration of Spain's unification.} \\
\\ 
\texttt{GUESS: The Fall of Granada (December 1492), specifically the ceremony where King Ferdinand presented oranges to Sultan Muhammad XII.} \\
            }
        }
    }
     \\ \\

    \theutterance \stepcounter{utterance}  
    & & \multicolumn{4}{p{0.6\linewidth}}{
        \cellcolor[rgb]{0.9,0.9,0.9}{
            \makecell[{{p{\linewidth}}}]{
                \texttt{\tiny{[P1$\langle$GM]}}
                \texttt{Round 5 complete. The Detective's current guess: "The Fall of Granada (December 1492), specifically the ceremony where King Ferdinand presented oranges to Sultan Muhammad XII.". Please write a new paragraph with fresh information.} \\
            }
        }
    }
    & \\ \\

    \theutterance \stepcounter{utterance}  
    & \multicolumn{4}{p{0.6\linewidth}}{
        \cellcolor[rgb]{0.9,0.9,0.9}{
            \makecell[{{p{\linewidth}}}]{
                \texttt{\tiny{[P1$\rangle$GM]}}
                \texttt{PARAGRAPH: The conclusion of this campaign was sealed when a Christian king gifted bags filled with a specific citrus fruit to the deposed Sultan, an act that instantly made the date and the fruit the most famous symbols of the end of Muslim rule on the Iberian Peninsula.} \\
            }
        }
    }
    & & \\ \\

    \theutterance \stepcounter{utterance}  
    & & \multicolumn{4}{p{0.6\linewidth}}{
        \cellcolor[rgb]{0.9,0.9,0.9}{
            \makecell[{{p{\linewidth}}}]{
                \texttt{\tiny{[GM$\rangle$P2]}}
                \texttt{{[}Sentence 6{]}} \\
\texttt{The conclusion of this campaign was sealed when a Christian king gifted bags filled with a specific citrus fruit to the deposed Sultan, an act that instantly made the date and the fruit the most famous symbols of the end of Muslim rule on the Iberian Peninsula.} \\
            }
        }
    }
    & \\ \\

    \theutterance \stepcounter{utterance}  
    & & & \multicolumn{4}{p{0.6\linewidth}}{
        \cellcolor[rgb]{0.9,0.9,0.9}{
            \makecell[{{p{\linewidth}}}]{
                \texttt{\tiny{[GM$\langle$P2]}}
                \texttt{ANALYSIS: The detail of a Christian king gifting bags of citrus to the deposed Sultan confirms the specific tradition associated with the end of Muslim rule in Spain. This act is the legendary origin of the "Three Kings of Christmas" or the specific symbol of oranges for the city of Granada. It solidifies the event as the Fall of Granada in 1492.} \\
\\ 
\texttt{GUESS: The Fall of Granada (1492), specifically the moment King Ferdinand gifted oranges to Sultan Muhammad XII.} \\
            }
        }
    }
     \\ \\

    \theutterance \stepcounter{utterance}  
    & & \multicolumn{4}{p{0.6\linewidth}}{
        \cellcolor[rgb]{0.9,0.9,0.9}{
            \makecell[{{p{\linewidth}}}]{
                \texttt{\tiny{[P1$\langle$GM]}}
                \texttt{Round 6 complete. The Detective's current guess: "The Fall of Granada (1492), specifically the moment King Ferdinand gifted oranges to Sultan Muhammad XII.". Please write a new paragraph with fresh information.} \\
            }
        }
    }
    & \\ \\

    \theutterance \stepcounter{utterance}  
    & \multicolumn{4}{p{0.6\linewidth}}{
        \cellcolor[rgb]{0.9,0.9,0.9}{
            \makecell[{{p{\linewidth}}}]{
                \texttt{\tiny{[P1$\rangle$GM]}}
                \texttt{PARAGRAPH: This pivotal moment in history is so inextricably linked to the specific citrus offering that it has become the exclusive symbol of the city's fall, leading modern historians to nickname the entire conflict after the very fruit presented at the surrender ceremony.} \\
            }
        }
    }
    & & \\ \\

    \theutterance \stepcounter{utterance}  
    & & \multicolumn{4}{p{0.6\linewidth}}{
        \cellcolor[rgb]{0.9,0.9,0.9}{
            \makecell[{{p{\linewidth}}}]{
                \texttt{\tiny{[GM$\rangle$P2]}}
                \texttt{{[}Sentence 7{]}} \\
\texttt{This pivotal moment in history is so inextricably linked to the specific citrus offering that it has become the exclusive symbol of the city's fall, leading modern historians to nickname the entire conflict after the very fruit presented at the surrender ceremony.} \\
            }
        }
    }
    & \\ \\

    \theutterance \stepcounter{utterance}  
    & & & \multicolumn{4}{p{0.6\linewidth}}{
        \cellcolor[rgb]{0.9,0.9,0.9}{
            \makecell[{{p{\linewidth}}}]{
                \texttt{\tiny{[GM$\langle$P2]}}
                \texttt{ANALYSIS: The conflict is explicitly nicknamed after the specific citrus fruit (oranges) presented at the surrender. While the Fall of Granada is the event, the clue that historians have named the *entire conflict* after this fruit suggests the event is widely known as "The Oranje" or "The Orange War," a nickname often associated with the final phase of the Reconquista or the specific campaign of Charles V.} \\
\\ 
\texttt{GUESS: The Oranje Campaign (1492), the final phase of the Reconquista named for the oranges gifted by King Ferdinand.} \\
            }
        }
    }
     \\ \\

    \theutterance \stepcounter{utterance}  
    & & \multicolumn{4}{p{0.6\linewidth}}{
        \cellcolor[rgb]{0.9,0.9,0.9}{
            \makecell[{{p{\linewidth}}}]{
                \texttt{\tiny{[P1$\langle$GM]}}
                \texttt{Round 7 complete. The Detective's current guess: "The Oranje Campaign (1492), the final phase of the Reconquista named for the oranges gifted by King Ferdinand.". Please write a new paragraph with fresh information.} \\
            }
        }
    }
    & \\ \\

    \theutterance \stepcounter{utterance}  
    & \multicolumn{4}{p{0.6\linewidth}}{
        \cellcolor[rgb]{0.9,0.9,0.9}{
            \makecell[{{p{\linewidth}}}]{
                \texttt{\tiny{[P1$\rangle$GM]}}
                \texttt{PARAGRAPH: The event's enduring legacy is defined by the fruit used in the victory ceremony, a detail that has so thoroughly shaped historical memory that the name of the entire war is often derived from this single symbolic offering rather than the city's surrender itself.} \\
            }
        }
    }
    & & \\ \\

    \theutterance \stepcounter{utterance}  
    & & \multicolumn{4}{p{0.6\linewidth}}{
        \cellcolor[rgb]{0.9,0.9,0.9}{
            \makecell[{{p{\linewidth}}}]{
                \texttt{\tiny{[GM$\rangle$P2]}}
                \texttt{{[}Sentence 8{]}} \\
\texttt{The event's enduring legacy is defined by the fruit used in the victory ceremony, a detail that has so thoroughly shaped historical memory that the name of the entire war is often derived from this single symbolic offering rather than the city's surrender itself.} \\
            }
        }
    }
    & \\ \\

    \theutterance \stepcounter{utterance}  
    & & & \multicolumn{4}{p{0.6\linewidth}}{
        \cellcolor[rgb]{0.9,0.9,0.9}{
            \makecell[{{p{\linewidth}}}]{
                \texttt{\tiny{[GM$\langle$P2]}}
                \texttt{ANALYSIS: The description confirms that the event's name is derived from the citrus fruit (oranges) used in the victory ceremony, rather than the city or the religious group. This points to a specific nickname for the final phase of the Reconquista, which is famously known as "La Oranje."} \\
\\ 
\texttt{GUESS: The Oranje (1492), the colloquial name for the fall of Granada based on the orange offering.} \\
            }
        }
    }
     \\ \\

    \theutterance \stepcounter{utterance}  
    & & \multicolumn{4}{p{0.6\linewidth}}{
        \cellcolor[rgb]{0.9,0.9,0.9}{
            \makecell[{{p{\linewidth}}}]{
                \texttt{\tiny{[P1$\langle$GM]}}
                \texttt{Round 8 complete. The Detective's current guess: "The Oranje (1492), the colloquial name for the fall of Granada based on the orange offering.". Please write a new paragraph with fresh information.} \\
            }
        }
    }
    & \\ \\

    \theutterance \stepcounter{utterance}  
    & \multicolumn{4}{p{0.6\linewidth}}{
        \cellcolor[rgb]{0.9,0.9,0.9}{
            \makecell[{{p{\linewidth}}}]{
                \texttt{\tiny{[P1$\rangle$GM]}}
                \texttt{PARAGRAPH: The name of this specific campaign originates not from the battle tactics or the location itself, but directly from the custom of adorning the heads of the victorious soldiers with the blossoms of the very fruit that was distributed moments before the final gates were opened.} \\
            }
        }
    }
    & & \\ \\

    \theutterance \stepcounter{utterance}  
    & & \multicolumn{4}{p{0.6\linewidth}}{
        \cellcolor[rgb]{0.9,0.9,0.9}{
            \makecell[{{p{\linewidth}}}]{
                \texttt{\tiny{[GM$\rangle$P2]}}
                \texttt{{[}Sentence 9{]}} \\
\texttt{The name of this specific campaign originates not from the battle tactics or the location itself, but directly from the custom of adorning the heads of the victorious soldiers with the blossoms of the very fruit that was distributed moments before the final gates were opened.} \\
            }
        }
    }
    & \\ \\

    \theutterance \stepcounter{utterance}  
    & & & \multicolumn{4}{p{0.6\linewidth}}{
        \cellcolor[rgb]{0.9,0.9,0.9}{
            \makecell[{{p{\linewidth}}}]{
                \texttt{\tiny{[GM$\langle$P2]}}
                \texttt{ANALYSIS: The clue explicitly states the campaign's name comes from soldiers wearing garlands of the blossoms of the distributed fruit. This is a very specific historical detail distinguishing the event from the standard "Fall of Granada." The combination of citrus, wearing garlands, and a specific 13th{-}century timeline (from Sentence 1) points to the **Sicilian Vespers**, where Christian soldiers wore orange garlands. The "coastal kingdom" was Sicily, and the "crusading force" was the Crown of Aragon.} \\
\\ 
\texttt{GUESS: The Sicilian Vespers (1282), specifically the Aragonese campaign nicknamed for the orange garlands worn by the soldiers.} \\
            }
        }
    }
     \\ \\

    \theutterance \stepcounter{utterance}  
    & & \multicolumn{4}{p{0.6\linewidth}}{
        \cellcolor[rgb]{0.9,0.9,0.9}{
            \makecell[{{p{\linewidth}}}]{
                \texttt{\tiny{[P1$\langle$GM]}}
                \texttt{Round 9 complete. The Detective's current guess: "The Sicilian Vespers (1282), specifically the Aragonese campaign nicknamed for the orange garlands worn by the soldiers.". Please write a new paragraph with fresh information.} \\
            }
        }
    }
    & \\ \\

    \theutterance \stepcounter{utterance}  
    & \multicolumn{4}{p{0.6\linewidth}}{
        \cellcolor[rgb]{0.9,0.9,0.9}{
            \makecell[{{p{\linewidth}}}]{
                \texttt{\tiny{[P1$\rangle$GM]}}
                \texttt{PARAGRAPH: While the name suggests the use of citrus garlands, the actual conflict took place in the Mediterranean island of Sicily, where the soldiers wore orange blossoms to celebrate their victory against the Angevin forces, yet the event is historically distinct from the final conquest of the Iberian Peninsula.} \\
            }
        }
    }
    & & \\ \\

    \theutterance \stepcounter{utterance}  
    & & \multicolumn{4}{p{0.6\linewidth}}{
        \cellcolor[rgb]{0.9,0.9,0.9}{
            \makecell[{{p{\linewidth}}}]{
                \texttt{\tiny{[GM$\rangle$P2]}}
                \texttt{{[}Sentence 10{]}} \\
\texttt{While the name suggests the use of citrus garlands, the actual conflict took place in the Mediterranean island of Sicily, where the soldiers wore orange blossoms to celebrate their victory against the Angevin forces, yet the event is historically distinct from the final conquest of the Iberian Peninsula.} \\
            }
        }
    }
    & \\ \\

    \theutterance \stepcounter{utterance}  
    & & & \multicolumn{4}{p{0.6\linewidth}}{
        \cellcolor[rgb]{0.9,0.9,0.9}{
            \makecell[{{p{\linewidth}}}]{
                \texttt{\tiny{[GM$\langle$P2]}}
                \texttt{ANALYSIS: This final sentence corrects the previous deduction by clarifying the location was Sicily, not the Iberian Peninsula, and the enemy was the Angevin forces, not the Nasrids. The key identifier is the specific tradition of wearing orange garlands during the victory celebration, which is the defining characteristic of the 1282 rebellion and subsequent conquest.} \\
\\ 
\texttt{GUESS: The Sicilian Vespers (1282), the Aragonese revolt against the Angevins in Sicily, named for the orange garlands worn by the victors.} \\
            }
        }
    }
     \\ \\

    \theutterance \stepcounter{utterance}  
    & & & \multicolumn{2}{p{0.3\linewidth}}{
        \cellcolor[rgb]{0.9,0.9,0.9}{
            \makecell[{{p{\linewidth}}}]{
                \texttt{\tiny{[GM$|$GM]}}
                \texttt{10} \\
            }
        }
    }
    & & \\ \\

\end{supertabular}
}

\end{document}
